\documentclass[11pt,a4paper]{article}
\usepackage[margin=25mm]{geometry}
\usepackage{fontspec}
\setmainfont{TeX Gyre Pagella}
\usepackage{amsmath,amssymb,mathtools}
\usepackage{graphicx}
\usepackage{xcolor}
\usepackage[unicode,colorlinks=true,linkcolor=blue,urlcolor=blue]{hyperref}
\usepackage{bookmark}
\usepackage{enumitem}
\setlist{itemsep=3pt,topsep=5pt}
\setlength{\emergencystretch}{3em}
\newcommand{\tightlist}{\setlength{\itemsep}{0pt}\setlength{\parskip}{0pt}}
\newcommand{\passthrough}[1]{#1}
\newcommand{\pandocbounded}[1]{#1}
\makeatletter
\def\maxwidth{\ifdim\Gin@nat@width>\linewidth\linewidth\else\Gin@nat@width\fi}
\def\maxheight{\ifdim\Gin@nat@height>0.75\textheight0.75\textheight\else\Gin@nat@height\fi}
\makeatother
\setkeys{Gin}{width=\maxwidth,height=\maxheight,keepaspectratio}
\hypersetup{pdfauthor={OpenAI GPT-6.1 Sol, Codex, Ultra effort}}
\begin{document}
\section{Power residue symbols and reciprocity
laws}\label{power-residue-symbols-and-reciprocity-laws}

\emph{Written by OpenAI GPT-6.1 Sol in Codex, Ultra effort, October
2026. Self-checked by the writing AI; no independent review is claimed.
Public domain (CC0).}

\textbf{Lesson 23.} A power residue symbol measures the Frobenius action
on a radical. The global product of local Hilbert symbols then compares
two power residue symbols with their arguments reversed. We recover
quadratic reciprocity and both supplements, prove a number-field primary
form of quadratic reciprocity, and compute the local factors needed for
cubic and quartic reciprocity. The primary normalizations are specified
separately for each law.

Use the Hilbert symbol of
\href{../hilbert-symbols-and-local-conics.html}{Hilbert symbols and
local conics}: the first argument enters arithmetic reciprocity and the
second is placed under the radical. Thus \[
(a,b)_{v,n}=\frac{\operatorname{rec}_{K_v,\mathrm{arith}}(a)(\sqrt[n]b)}{\sqrt[n]b}.
\tag{1}
\] Proposition 11.1 proves multiplicativity, the norm criterion and
\((a,b)_{v,n}=(b,a)_{v,n}^{-1}\); Proposition 11.3 proves the tame
formula. \href{../the-global-reciprocity-law.html}{The global
reciprocity law}, especially its principal product and local
compatibility, supplies the global input. No power reciprocity law is
assumed in those proofs.

\subsection{1. Residue symbols and their Frobenius
interpretation}\label{residue-symbols-and-their-frobenius-interpretation}

Let \(K\) be a global field containing \(\mu_n\), with \(n\geq2\) prime
to its characteristic. At a finite prime \(\mathfrak p\nmid n\), put
\(q=N\mathfrak p\), and suppose \(a\in K^\times\) is a local unit there.
Reduction embeds \(\mu_n\) into \(\kappa(\mathfrak p)^\times\): a
nontrivial \(n\)-th root cannot be congruent to 1, by the
derivative-unit argument of lesson 11. In particular \(n\mid q-1\).
Define the \textbf{power residue symbol} by the unique root \[
\left(\frac a{\mathfrak p}\right)_n\in\mu_n,
\qquad
\left(\frac a{\mathfrak p}\right)_n\equiv a^{(q-1)/n}\pmod{\mathfrak p}.
\tag{2}
\] The displayed residue is an \(n\)-th root, by the finite-field
multiplicative group order.

\textbf{Proposition 23.1.} The symbol depends only on the nonzero
residue of \(a\), is multiplicative in \(a\), and equals 1 exactly when
that residue is an \(n\)-th power. For a fractional ideal
\(\mathfrak b=\prod\mathfrak p^{m_{\mathfrak p}}\), supported away from
\(n\) and at places where \(a\) is a unit, the definition \[
\left(\frac a{\mathfrak b}\right)_n
=\prod_{\mathfrak p}\left(\frac a{\mathfrak p}\right)_n^{m_{\mathfrak p}}
\tag{3}
\] is multiplicative in both its numerator and ideal denominator.

\textbf{Proof.} Uniqueness in (2) gives residue dependence and
multiplicativity. The finite-field group is cyclic. One elementary proof
is to let \(m\) be the exponent of this finite abelian group: for every
prime power dividing \(m\), an element attaining that prime-power order
exists; their product has order \(m\). Every group element is a root of
\(X^m-1\), so the root bound gives group order at most \(m\), and hence
equal to \(m\). Thus it has a generator. If \(g\) is a generator of
order \(q-1\), then \((g^r)^{(q-1)/n}=1\) exactly when \(n\mid r\),
exactly when \(g^r\) is an \(n\)-th power. Formula (3) and its
multiplicativity follow from unique prime ideal factorization, with
negative exponents permitted for fractional ideals. \(\square\)

For a principal denominator write \((a/b)_n=(a/(b))_n\). A composite
denominator loses information: even for \(n=2\), a Jacobi symbol of 1
need not mean a square modulo that denominator. For example
\((2/15)_2=(2/3)_2(2/5)_2=1\), whereas 2 is a square modulo neither 3
nor 5.

\textbf{Proposition 23.2.} Let \(\beta^n=a\). The extension
\(K(\beta)/K\) is unramified at \(\mathfrak p\) as above, and its
arithmetic Frobenius satisfies \[
\frac{\operatorname{Frob}_{\mathfrak p}(\beta)}\beta
=\left(\frac a{\mathfrak p}\right)_n
=(\pi_{\mathfrak p},a)_{\mathfrak p,n}
=(a,\pi_{\mathfrak p})_{\mathfrak p,n}^{-1}.
\tag{4}
\] The uniformizer can be any local prime element.

\textbf{Proof.} In a finite residue extension containing an \(n\)-th
root of \(\bar a\), the derivative of \(X^n-a\) at that unit root is a
unit. The unramified lifting theorem and Hensel's iteration give a root
in the corresponding unramified local extension. Hence the completion of
\(K(\beta)\) is unramified. Its Frobenius ratio reduces to
\(\bar\beta^{q-1}=\bar a^{(q-1)/n}\). That ratio is in \(\mu_n\), so
reduction injectivity gives the first equality of (4). Arithmetic local
reciprocity sends a uniformizer to precisely this Frobenius, giving the
second equality through (1). Skew symmetry gives the inverse in the last
equality. A change of uniformizer is a unit, which acts trivially on the
unramified radical extension. The argument does not assume its degree
equals \(n\). \(\square\)

\subsection{2. The global product and general
reciprocity}\label{the-global-product-and-general-reciprocity}

\textbf{Theorem 23.3 (Hilbert's product formula).} For
\(a,b\in K^\times\), \[
\prod_v(a,b)_{v,n}=1.
\tag{5}
\] Only finitely many factors differ from 1.

\textbf{Proof.} The extension \(L=K(\sqrt[n]b)\) is finite cyclic, since
\(\mu_n\subset K\); its degree may divide \(n\). At almost every finite
place, \(a\) and \(b\) are units and \(n\) is a unit, so the radical
extension is unramified and reciprocity of \(a\) is trivial. For the
finitely many other places, global local compatibility says that the
product of the local reciprocity automorphisms is the global symbol of
the principal idèle \(a\). This is identity by global reciprocity.
Evaluate it on a fixed global root of \(b\), using compatible embeddings
in each completion. Its ratio is the product in (5), so that product is
1. \(\square\)

At a tame finite place write \(a=\pi^r u\), \(b=\pi^t z\), with units
\(u,z\). The precise tame formula of lesson 11 is \[
\overline{(a,b)_{v,n}}
=\left((-1)^{rt}\frac{\bar z^{\,r}}{\bar u^{\,t}}\right)^{(Nv-1)/n}.
\tag{6}
\] In particular two units pair trivially.

\textbf{Theorem 23.4 (general power reciprocity).} Suppose
\(a,b\in\mathcal O_K\) are nonzero, their principal ideals are coprime,
and both are prime to \(n\). Then \[
\left(\frac a b\right)_n
\left(\frac b a\right)_n^{-1}
=\prod_{v\mid n\infty}(a,b)_{v,n}.
\tag{7}
\] For function fields, the same statement uses the chosen valuation
supports in place of \(\mathcal O_K\); outside the indicated exceptional
places they must be disjoint. It also holds for fractional principal
ideals with disjoint supports and unit components at the exceptional
finite places. In positive characteristic prime to \(n\), there are no
places dividing \(n\) and no infinite places, so the right side is empty
and equals 1.

\textbf{Proof.} At a prime of \((b)\), the numerator \(a\) is a unit and
(4) and multiplicativity give \((a,b)_{v,n}=(a/v)_n^{-v(b)}\). At a
prime of \((a)\), similarly \((a,b)_{v,n}=(b/v)_n^{v(a)}\). At all other
tame places both are units, and (6) gives 1. Separate these factors in
(5) to obtain \[
1=\left(\frac b a\right)_n
  \left(\frac a b\right)_n^{-1}
  \prod_{v\mid n\infty}(a,b)_{v,n},
\] which is (7). The same argument works with negative valuations for
fractional ideals. \(\square\)

Thus a Frobenius residue symbol is \((\pi,a)\), whereas reversing the
entries to \((a,\pi)\) inserts an inverse. For quadratic signs this
inverse is invisible; for higher powers it is essential.

\subsection{3. Quadratic reciprocity over the
rationals}\label{quadratic-reciprocity-over-the-rationals}

For distinct odd positive primes \(p,q\), (7) with \(n=2\) has only the
2-adic and real exceptional factors. The real factor is 1 because both
entries are positive. The explicit 2-adic unit formula proved in lesson
11 gives \[
(p,q)_{2,2}=(-1)^{((p-1)/2)((q-1)/2)}.
\tag{8}
\]

\textbf{Corollary 23.5.} The Legendre symbols satisfy \[
\left(\frac p q\right)\left(\frac q p\right)
=(-1)^{((p-1)/2)((q-1)/2)},\quad
\left(\frac{-1}p\right)=(-1)^{(p-1)/2},\quad
\left(\frac2p\right)=(-1)^{(p^2-1)/8}.
\tag{9}
\]

\textbf{Proof.} For the first equation, inversion equals identity on
signs, so (7) and (8) give the assertion. For \((-1/p)\), use (5) on
\((-1,p)\). At \(p\), (4) gives the inverse of \((-1/p)\), hence the
same sign. The real factor is 1 because \(p>0\). Other odd places have
two unit entries and factor 1. The 2-adic factor, by the unit formula,
is \((-1)^{(p-1)/2}\). Their product is 1, giving the first supplement.

For \((2/p)\), use (5) on \((2,p)\). The only possible nontrivial
factors are at 2 and \(p\); the real factor is again 1. At \(p\) the
factor is \((2/p)^{-1}=(2/p)\). The full 2-adic formula, with valuation
of the first entry 1 and the second entry an odd unit, gives
\((2,p)_{2,2}=(-1)^{(p^2-1)/8}\). The product proves the second
supplement. \(\square\)

For example 2 is a square modulo an odd prime exactly when the prime is
1 or 7 modulo 8. This conclusion comes from the local obstruction at 2,
rather than from an unexplained global sign.

\subsection{4. A primary quadratic law over any number
field}\label{a-primary-quadratic-law-over-any-number-field}

In this section a number \(\alpha\in\mathcal O_K\) is \textbf{primary
for the quadratic law} if it is odd and
\(\alpha\equiv\xi^2\pmod{4\mathcal O_K}\) for some
\(\xi\in\mathcal O_K\). This definition belongs to this quadratic
statement; the cubic and quartic normalizations below are different.

For odd, coprime \(\alpha,\beta\), with at least one primary in this
sense, the general quadratic law takes the particularly simple form \[
\left(\frac\alpha\beta\right)_2
\left(\frac\beta\alpha\right)_2
=(-1)^{\#\{v\text{ real}:\alpha_v<0,\ \beta_v<0\}}.
\tag{10}
\] In particular it is symmetric if at least one of the entries is
totally positive. This is the primary quadratic reciprocity law in the
notation of our symbols.

\textbf{Proof.} Suppose \(\alpha\) is primary. At each dyadic completion
\(F\), \(\xi\) is a unit and \(\alpha/\xi^2=1+4c\),
\(c\in\mathcal O_F\). An element \(\theta\) satisfying
\(\theta^2+\theta=c\) gives \((1+2\theta)^2=1+4c\). Its equation reduces
to a separable quadratic over the finite residue field, since its
derivative \(2\theta+1\) is a unit. If the reduced equation has a root
it lifts in \(F\); otherwise it lifts in the unramified quadratic
extension. Hence \(F(\sqrt\alpha)/F\) is trivial or unramified. The unit
\(\beta\) has trivial reciprocity on it, so \((\beta,\alpha)_{F,2}=1\),
and skew symmetry gives \((\alpha,\beta)_{F,2}=1\). Thus every dyadic
factor in (7) is 1. Complex factors are 1. At a real place the quadratic
Hilbert symbol is \(-1\) exactly when both entries are negative, as
proved in lesson 11. Their product is (10). The argument with
\(\alpha,\beta\) exchanged is identical. \(\square\)

The proof above uses the unramified local quadratic equation and the
already proved global product, with no hypothesis that the field's
different is principal. The more general law without a primary entry
remains (7), retaining its dyadic factors.

\subsection{5. Cubic reciprocity and its local
calculation}\label{cubic-reciprocity-and-its-local-calculation}

Put \(\omega=e^{2\pi i/3}\), \(K=\mathbf Q(\omega)\), and
\(\lambda=1-\omega\). Its integer ring is \(\mathbf Z[\omega]\);
\(\lambda^2=-3\omega\), and the unique prime over 3 has residue field
\(\mathbf F_3\). We call an element prime to 3 \textbf{primary for the
cubic law} if it is congruent to \(-1\pmod3\). Each prime away from 3
has a unique associate with this normalization. Indeed the six units
\(\{\pm1,\pm\omega,\pm\omega^2\}\) have distinct residues and exhaust
the six units modulo \(\lambda^2\); multiplication by exactly one gives
residue \(-1\). Their distinctness follows because differences between
distinct cube roots have valuation 1, and changing sign gives a unit
difference. The norm \(a^2-ab+b^2\) also shows directly that these are
all the global units. Allowing either sign modulo 3 gives the alternate
two-associate convention; \(-1\) is a cube, so it changes no cubic
symbol.

The necessary wild local assertion will be proved explicitly. Let
\(F=\mathbf Q_3(\omega)\), with \(U^j=1+\lambda^j\mathcal O_F\).

\textbf{Cubic local lemma.} If \(a,b\in U^2\), then \((a,b)_{F,3}=1\).

\textbf{Proof.} Write \(b=1+\lambda^2 c\). If \(b\) is a cube, the
assertion is immediate. Otherwise \(L=F(\gamma)\), \(\gamma^3=b\), is
cyclic of degree three. For \(t\in\mathcal O_F\), \[
N_{L/F}(1+t(\gamma-1))=(1-t)^3+bt^3
=1+\lambda^2 P(t),
\quad P(t)=\omega^2t-\omega^2t^2+ct^3.
\tag{11}
\] The first identity is the product over the three conjugates of
\(\gamma\); the second uses \(3=-\omega^2\lambda^2\). Since
\(P'(0)=\omega^2\) is a unit, for each \(d\in\lambda\mathcal O_F\)
Hensel's iteration at zero solves \(P(t)=d\) with
\(t\in\lambda\mathcal O_F\). Thus every element of \(U^3\) is a norm by
(11).

Modulo \(\lambda\), \(P(1)=c\) and \(P(-1)=1-c\). At least one of these
residues is nonzero in \(\mathbf F_3\). The corresponding norm in (11)
lies in \(U^2\) and has nonzero image in
\(U^2/U^3\simeq(\mathbf F_3,+)\), so its image generates that quotient.
The norm group already contains all of \(U^3\), hence contains all of
\(U^2\): adjust any target's residue by a power of that norm and its
remaining part is in \(U^3\). Every element used in (11) is nonzero,
since its norm is a unit. The norm criterion of Proposition 11.1 now
gives \((a,b)_{F,3}=1\). \(\square\)

\textbf{Cubic reciprocity.} For nonassociate primary primes \(\pi,\rho\)
away from 3, \[
\left(\frac\pi\rho\right)_3=\left(\frac\rho\pi\right)_3.
\tag{12}
\] The same holds for coprime primary elements.

\textbf{Proof.} Their negatives lie in \(U^2\) at \(\lambda\), and their
signs are cubes, so the lemma makes the local exceptional factor 1. All
infinite places are complex. Therefore (7) has right side 1 and gives
(12). Multiplicativity gives the version for coprime primary elements
directly from the same local argument. \(\square\)

The use of the full residue cardinality matters. For a rational prime
remaining prime in \(\mathbf Z[\omega]\), that cardinality is \(p^2\),
not \(p\). Formula (2) works without distinguishing split and inert
primes. No unspecified local cubic symbol is left in the proof.

\subsection{6. Quartic reciprocity and a two-generator local
pairing}\label{quartic-reciprocity-and-a-two-generator-local-pairing}

Now put \(K=\mathbf Q(i)\), \(\lambda=1+i\). An odd Gaussian element is
\textbf{primary for the quartic law} if it is 1 modulo \(\lambda^3\),
equivalently modulo \(2+2i\), since these generators are associates. The
four units \(\mu_4\) give all four distinct unit residues modulo
\(\lambda^3\), so every odd Gaussian prime has a unique primary
associate. For \(a+bi\) this congruence says that \(a\) is odd, \(b\) is
even, and \(a+b\equiv1\pmod4\).

Work locally in \(F=\mathbf Q_2(i)\), with
\(U^j=1+\lambda^j\mathcal O_F\). The integer ring is \(\mathbf Z_2[i]\),
as follows from the Eisenstein equation for \(\lambda\). For
\(u\in U^3\), its norm is 1 modulo 4, by the coordinate description just
given. Define \[
e(u)=\frac{N_{F/\mathbf Q_2}(u)-1}{4}\pmod2.
\tag{13}
\] It is a homomorphism on \(U^3\): multiply two norms congruent to 1
modulo 4 and reduce the quotient by 4 modulo 2.

\textbf{Quartic local lemma.} For \(u,v\in U^3\), \[
(u,v)_{F,4}=(-1)^{e(u)e(v)}.
\tag{14}
\]

\textbf{Proof.} First \((U^3)^2=U^5\). If \(x\in\lambda^3\mathcal O_F\),
then \((1+x)^2-1=2x+x^2\in\lambda^5\mathcal O_F\). Conversely for
\(y\in\lambda^5\mathcal O_F\), solve \(x=y/2-x^2/2\) by contraction in
\(\lambda^3\mathcal O_F\). The constant term has valuation at least 3,
the quadratic term at least 4, and the difference at two arguments is
multiplied by \((x+z)/2\), of valuation at least 1. The iterates
converge to the required square root in \(U^3\).

The quotient \(U^3/U^5\) has order four, since both successive residue
quotients have order two. It has basis \[
g=-1+2i,\qquad h=5.
\tag{15}
\] Indeed \(v_F(g-1)=3\), \(v_F(h-1)=4\), and \(g^2-1=-4-4i\),
\(h^2-1=24\) have valuations 5 and 6. Thus both have order two in the
quotient and are independent by their distinct leading levels.

We compute three quartic pairings using the already proved global
product formula over \(\mathbf Q(i)\) and the tame formula. This uses no
quartic reciprocity theorem. The only odd places involved lie over 5,
generated by \(g\) and \(\bar g=-1-2i\); both have norm 5, and
\(5=g\bar g\). For \((g,g)_4\), the tame factor at \(g\) is \(-1\), and
every other odd factor is 1. For \((h,h)_4\), each of the two places
over 5 has tame factor \(-1\), with product 1. All infinite factors are
trivial. Hence (5) gives \[
(g,g)_{F,4}=-1,\qquad(h,h)_{F,4}=1.
\tag{16}
\] For \((g,h)_4\), at \(g\) formula (6) gives the residue
\(-\bar g=2\pmod g\). Here \(i=3\pmod g\), so that factor is \(-i\). At
\(\bar g\), it gives \(g^{-1}=2\pmod{\bar g}\), and
\(i=2\pmod{\bar g}\), so that factor is \(i\). Their product is 1,
yielding \[
(g,h)_{F,4}=1.
\tag{17}
\]

Compatibility of radicals gives \((a,b)_4^2=(a,b)_2\): the square of a
fourth root of \(b\) is a square root. The squares of all pairings in
(16)--(17) are 1, and the quadratic symbol is skew symmetric. Since
\(U^5=(U^3)^2\), its square classes are already accounted for by the two
generators (15). Bilinearity therefore shows that the quadratic pairing
is trivial on all of \(U^3\times U^3\). Consequently all quartic values
there are signs. Also \(U^5\) is in the radical of that restricted
quartic pairing, since for \(z,v\in U^3\),
\((z^2,v)_4=(z,v)_4^2=(z,v)_2=1\), and likewise in the other argument.

Thus (16)--(17) describe the entire restricted pairing on the
two-dimensional binary quotient \(U^3/U^5\): its only negative basis
pairing is the \(g,g\) entry. Finally \(e(g)=1\), \(e(h)=0\), and \(e\)
kills \(U^5\), since the square of a norm congruent to 1 modulo 4 is 1
modulo 8. Its value is precisely the coefficient of \(g\) in this
quotient. This proves (14). \(\square\)

\textbf{Quartic reciprocity.} For distinct primary Gaussian primes
\(\pi,\rho\), \[
\left(\frac\pi\rho\right)_4
=\left(\frac\rho\pi\right)_4
(-1)^{((N\pi-1)/4)((N\rho-1)/4)}.
\tag{18}
\]

\textbf{Proof.} Both local entries at the unique dyadic prime lie in
\(U^3\). Their local norms are their ordinary Gaussian norms, so
(13)--(14) give exactly the sign in (18). This is the only exceptional
finite factor in (7), and the infinite places are complex. Substitute it
into (7) to obtain (18). \(\square\)

The term biquadratic reciprocity also denotes this fourth-power law. Its
sign is absent from cubic reciprocity because the corresponding primary
cubic local pairing vanishes. The local unit quotient computation above
explains the quartic sign rather than imposing it as a convention.

\subsection{7. Exercises and complete
solutions}\label{exercises-and-complete-solutions}

\subsubsection{Exercise 1 --- The second quadratic supplement
(easy)}\label{exercise-1-the-second-quadratic-supplement-easy}

Derive \((2/p)=(-1)^{(p^2-1)/8}\) for an odd positive prime from the
product formula.

\textbf{Solution.} In \(\prod_v(2,p)_{v,2}=1\), every odd place other
than \(p\) has two unit entries and tame factor 1. The real factor is 1.
The factor at \(p\) is the inverse of \((2/p)\) by (4), which is the
same sign. The explicit dyadic formula gives
\((2,p)_{2,2}=(-1)^{(p^2-1)/8}\). Thus those two remaining signs
multiply to 1 and are equal. The exponent is even for
\(p\equiv1,7\pmod8\) and odd for \(p\equiv3,5\pmod8\), giving the
familiar residue criterion.

\subsubsection{Exercise 2 --- Recovering the symbol from Frobenius
(medium)}\label{exercise-2-recovering-the-symbol-from-frobenius-medium}

Show that \((a/\mathfrak p)_n\) is determined by the Frobenius in
\(K(\sqrt[n]a)/K\), including when the radical extension has degree
smaller than \(n\).

\textbf{Solution.} Suppose \(a\) is a unit and \(\mathfrak p\nmid n\). A
residue root lifts in an unramified local extension because the
derivative is a unit. Choose a global root \(\beta\) and a place above
\(\mathfrak p\). Arithmetic Frobenius sends its residue to
\(\bar\beta^q\), so its ratio with \(\beta\) is a root in \(\mu_n\)
reducing to \(\bar a^{(q-1)/n}\). Reduction injectivity on \(\mu_n\)
makes this ratio exactly (2). Replacing \(\beta\) by another root
multiplies it by an element of \(\mu_n\subset K\), fixed by Frobenius,
so changes no ratio. Conjugating the chosen place changes no action in
this abelian radical extension. Nothing in these steps asserts degree
\(n\); the Frobenius may have smaller order and its ratio then lies in
the corresponding subgroup of \(\mu_n\). Finally it is \((\pi,a)_n\),
and equals the inverse of \((a,\pi)_n\) with the convention (1).

\subsubsection{\texorpdfstring{Exercise 3 --- The quartic symbol of
\(i\)
(medium)}{Exercise 3 --- The quartic symbol of i (medium)}}\label{exercise-3-the-quartic-symbol-of-i-medium}

For an odd Gaussian prime \(\pi\), compute \((i/\pi)_4\). Give examples
with values \(i\) and \(-1\).

\textbf{Solution.} The residue cardinality satisfies \(4\mid N\pi-1\),
since \(\mu_4\) injects into its multiplicative group. Formula (2) says
the desired root reduces to \(i^{(N\pi-1)/4}\). This is already a global
fourth root of unity, so injectivity gives the exact equality \[
\left(\frac i\pi\right)_4=i^{(N\pi-1)/4}.
\tag{19}
\] For \(\pi=-1+2i\), norm 5, the exponent is 1 and the value is \(i\).
For the inert Gaussian prime \(\pi=-3\), norm 9, it is 2 and the value
is \(-1\). Both chosen associates are primary, although (19) requires no
primary normalization. Using norm 3 instead of 9 for the second prime
would not even give an integer exponent and would be incorrect.

\subsubsection{Exercise 4 --- The primary cubic local norm calculation
(hard)}\label{exercise-4-the-primary-cubic-local-norm-calculation-hard}

Prove cubic reciprocity for nonassociate primary primes by computing the
exceptional Hilbert symbol at \(1-\omega\).

\textbf{Solution.} Primary primes are \(-1\pmod3\). Their negatives are
in \(U^2\subset\mathbf Q_3(\omega)^\times\), and \(-1\) is a cube. It is
enough to prove every element of \(U^2\) is a norm from the radical
field of any \(b\in U^2\). If \(b\) is a cube this is automatic.
Otherwise write \(b=1+\lambda^2c\) and \(\gamma^3=b\). Compute the norm
of \(1+t(\gamma-1)\) as in (11). The derivative \(P'(0)=\omega^2\) is a
unit, so for every \(d\in\lambda\mathcal O_F\) the equation \(P(t)=d\)
has a root near zero. Consequently all of \(U^3\) is in the norm group.
Among \(P(1)\) and \(P(-1)\), the residues are \(c\) and \(1-c\); at
least one is nonzero and its norm generates the three-element quotient
\(U^2/U^3\). The norm group therefore contains all of \(U^2\). The norm
criterion proves that the Hilbert symbol of the two primary entries is 1
at \(\lambda\). This is the only exceptional finite place, and complex
symbols are trivial. The general reciprocity equation (7) gives (12).
This supplies the local calculation itself, rather than citing a
wild-symbol formula without its proof.

\subsection{What this lesson does not
prove}\label{what-this-lesson-does-not-prove}

The Kummer, local Hilbert-symbol and global principal reciprocity
results are the exact earlier written course inputs specified at the
beginning. All five numbered results, the arbitrary-number-field primary
quadratic theorem, and the cubic and quartic primary laws are proved
here. The necessary wild primary local calculations are given explicitly
in sections 5--6. Further unit and ramified-prime supplementary laws for
higher powers are not asserted here.

\subsection{Editable edition}\label{editable-edition}

The reading edition provides the complete LaTeX source of this lesson,
the cumulative course LaTeX and the editable source ZIP. The archive
contains all twenty-four Markdown lessons, complete LaTeX bodies,
original diagrams, metadata and reproduction instructions.

\subsection{References}\label{references}

The global Hilbert product and the tame formula give the general
power-residue identity. The primary quadratic law keeps its real sign.
The explicit cubic norm polynomial and quartic local two-generator
pairing prove the primary cubic and quartic laws in the lesson.

\begin{itemize}
\tightlist
\item
  \href{https://www.jmilne.org/math/CourseNotes/CFT.pdf}{J. S. Milne,
  Class Field Theory, version 4.03}.
\end{itemize}

The \href{../free-proofs.html}{proof guide} gives the lesson sequence
and the exact prerequisite record. External references accompany the
written arguments.

\end{document}
